\chapter{Results}\label{chap:results}

\section{Full-Reference Fusion Metric Performance}

Following the SVD and Picard analysis from Fig.~\ref{fig:svd_picard}, which indicated that $k=6$ reached 95\% explained variance with stable projections, Table~\ref{tab:fr_individual_metrics} reports the marginal correlation and error of the selected FR metrics with MOS.

\begin{table}
\centering
\caption{Individual FR-IQA metrics versus MOS on the unbias track. Higher PLCC and SRCC indicate stronger correlation; lower MSE and MAE indicate more accurate estimates.}
\label{tab:fr_individual_metrics}
\begin{tabular}{lcccc}
\hline
Metric & PLCC & SRCC & MSE & MAE \\
\hline
MS-SSIM & 0.7712 & \textbf{0.8489} & 2542.00 & 47.27 \\
PSNR & \textbf{0.8009} & 0.8132 & \textbf{420.70} & \textbf{16.57} \\
PSNR-B & 0.7996 & 0.8126 & 434.80 & 16.84 \\
SNR & 0.7931 & 0.8047 & 493.90 & 18.08 \\
VIF & 0.7723 & 0.7670 & 2584.00 & 47.75 \\
UQI & 0.6865 & 0.8085 & 2540.00 & 47.24 \\
\hline
\end{tabular}
\end{table}

Each regressor is tuned with a TPE-based hyperparameter search in Optuna with a Successive Halving pruner and evaluated by 5-fold cross-validation. The selection metric is the mean of PLCC and SRCC between cross-validated predictions and MOS. Table~\ref{tab:fusion_models_k6} compares all fusion regressors trained at $k=6$.

\begin{table}
\centering
\caption{Fusion regressors at $k=6$ on the unbias track.}
\label{tab:fusion_models_k6}
\begin{tabular}{lcccc}
\hline
Model & PLCC & SRCC & MSE & MAE \\
\hline
CatBoost & \textbf{0.8932} & \textbf{0.8950} & \textbf{62.28} & \textbf{6.046} \\
GradientBoosting & 0.8901 & 0.8905 & 63.98 & 6.100 \\
XGBoost & 0.8891 & 0.8915 & 64.55 & 6.157 \\
ExtraTrees & 0.8871 & 0.8892 & 65.65 & 6.131 \\
RandomForest & 0.8861 & 0.8879 & 66.25 & 6.094 \\
LightGBM & 0.8838 & 0.8858 & 67.66 & 6.313 \\
HistGradientBoosting & 0.8760 & 0.8790 & 71.94 & 6.507 \\
Ridge & 0.8131 & 0.8353 & 105.71 & 7.941 \\
ElasticNet & 0.8129 & 0.8354 & 106.55 & 7.983 \\
BayesianRidge & 0.8145 & 0.8302 & 103.64 & 7.761 \\
Linear & 0.8147 & 0.8284 & 103.55 & 7.754 \\
\hline
\end{tabular}
\end{table}

The best fusion model is CatBoost configured with 3{,}000 iterations, a learning rate of 0.0118, depth 10, and an $\ell_2$ leaf-regularization coefficient of 1.005, selected via Optuna. Post hoc, we compute SHAP values to interpret the fitted CatBoost model. As shown in Fig.~\ref{fig:shap_summary}, MS-SSIM and UQI exert the strongest influence on the fused prediction, while VIF, PSNR-B, SNR, and PSNR have moderate effects.


\begin{figure}
    \centering
    \includegraphics[width=0.9\linewidth]{images/shap_summary.png}
    \caption{SHAP summary plot for $k=6$, computed on the best CatBoost fusion model.}
    \label{fig:shap_summary}
\end{figure}

\section{No-Reference Regression Model Performance}

We refer to our NR-IQA model as pMOS-Face (pseudo-MOS for Face), reflecting its use of pseudo-MOS labels generated from the FR fusion model and its evaluation on faces. This framework can be adapted to other datasets, provided that a small set of images with human labels is available for training the FR fusion model. Fig.~\ref{fig:nr-iqa-metric} shows two evaluation scenarios: in Fig.~\ref{fig:nr-iqa-metric}(a), the model's predictions are compared against the pseudo and ground-truth MOS used during training, yielding a PLCC of 0.9804, SRCC of 0.9770, MSE of 6.165, and MAE of 1.638. Fig.~\ref{fig:nr-iqa-metric}(b) shows the predictions compared against the MOS test set, resulting in a PLCC of 0.9679, SRCC of 0.9691 and MSE of 35.34 and MAE of 4.743. These results confirm both generalization to unseen human labels and alignment with the pseudo-supervision used for training. % chktex 36

\begin{figure}[ht]
    \centering
    \begin{minipage}[t]{0.48\textwidth}
        \centering
        \includegraphics[width=\linewidth]{images/predictions_train.pdf}
        \textbf{(a)} Training plot.\@
    \end{minipage}
        \hfill
        \begin{minipage}[t]{0.48\textwidth}
        \centering
        \includegraphics[width=\linewidth]{images/predictions_test.pdf}
        \textbf{(b)} Test plot.\@
    \end{minipage}
    \caption{Evaluation of our NR-IQA model. Panel (a) shows performance on the training set with pseudo and ground-truth MOS;\@ panel (b) shows performance on the disjoint test set with ground-truth MOS.}\label{fig:nr-iqa-metric}
\end{figure}

The results indicate that our model is more effective in predicting perceptual quality than classical NR-IQA metrics such as NIQE~\cite{mittal2012niqe} and PIQE~\cite{venka2015piqe}, as well as task-specific approaches like FaceQnet~\cite{hernandez2019faceqnet}, SER-FIQ~\cite{terhorst2020serfiq}, and MagFace~\cite{meng2021magface}. This highlights the robustness of our weakly supervised strategy for NR-IQA on steganographically degraded facial images. Quantitative results are reported in Table~\ref{tab:nr_scores}.

\begin{table}
\centering
\caption{Performance of our pMOS-Face compared to standard NR-IQA baselines.}\label{tab:nr_scores}
\begin{tabular}{lcccc}
\hline % chktex 44
\textbf{Metric} & \textbf{PLCC} & \textbf{SRCC} & \textbf{MSE} & \textbf{MAE} \\
\hline % chktex 44
pMOS-Face (ours) & \textbf{0.9285} & \textbf{0.9345} & \textbf{120.6} & \textbf{9.158} \\
NIQE                 & 0.7536          & 0.7431          & 6859        & 82.62 \\
PIQE                 & 0.2993          & 0.3114          & 3297        & 57.05 \\
BRISQUE              & 0.1234          & 0.5678          & 9101        & 23.45 \\
FaceQnet             & 0.4567          & 0.4789          & 5123        & 34.12 \\
SER-FIQ              & -0.1648         & -0.1542         & 123.5         & 9.230 \\
MagFace              & -0.6095         & -0.6362         & 4437        & 66.24 \\
\hline % chktex 44
\end{tabular}
\end{table}

When considering a more generalized use, where observer bias is not accounted for, our NR-IQA model shows a major improvement over the baselines. This is particularly relevant for task specific applications, where the model can be trained on a small set of images with human labels and then applied to larger datasets without the need for additional supervision. This approach allows for a more scalable and efficient solution not only for FIQA but also for other domains where reference images are not available.
